\documentclass[11pt]{amsart}
\usepackage[T1]{fontenc}
\usepackage{lmodern}
\usepackage{amsmath,amssymb,amsthm}
\usepackage[numbers,sort]{natbib}
\usepackage[hidelinks]{hyperref}
\usepackage{bm}

\newtheorem{theorem}{Theorem}
\newtheorem{lemma}[theorem]{Lemma}
\newtheorem{corollary}[theorem]{Corollary}
\theoremstyle{definition}
\newtheorem{definition}[theorem]{Definition}
\newtheorem{remark}[theorem]{Remark}

\newcommand{\Ter}{\mathcal S}
\newcommand{\red}{\twoheadrightarrow}
\newcommand{\hred}{\twoheadrightarrow_h}
\newcommand{\hstep}{\to_h}
\newcommand{\conv}{=_\beta}
\newcommand{\sz}[1]{\lvert#1\rvert}
\newcommand{\NF}{\mathrm{NF}}
\newcommand{\rest}{\mathrm{rest}}
\newcommand{\Sk}[1]{S_{\times#1}}

\title[Determined by Their Trees]{A Family of $\bm{S}$-Terms Determined by Their Infinite Trees}
\author{Claude Opus \& Fable with N.\ E.\ Davis}
\address{ZeroAttest, Champaign, Illinois, USA}
\email{NDavis@ZeroAttest.com}
\urladdr{https://orcid.org/0000-0002-1131-8319}
\date{\today}

\begin{document}
\begin{abstract}
Every closed term built from the combinator $S$ has an infinitary normal form, its tree, and
convertible terms have the same tree. Whether terms with the same tree are always convertible is
open. Let $e=S\,S$, $Q=S\,e\,e$ and $G_i=e\,(e\,(\cdots(e\,Q)\cdots))$ with $i$ copies of $e$. We
prove that if $Y\,G_m$ has the same tree as $G_i\,G_m$, where $i\ge0$ and $m\ge1$, then $Y$ is
$G_i$ as written. The 38-element block model of the head-normal-form theorem confines $Y$ to three
blocks, a finite fact about its table fixes the shape of the head normal form of $Y$, and an
induction on $i$ does the rest. As a corollary, every closed $S$-term with the same tree as
$G_i\,G_m$ is convertible with $G_i\,G_m$. These terms do not normalize, and every node on the
rightmost path of their trees has a head normal form. The proof rests on finitely many computed
facts about the table of the block model and on the invariance of the block under tree equality,
proved in a companion paper. The theorem and the corollary are mechanized in Lean~4, with the
table facts checked by kernel evaluation, using only the axioms of propositional extensionality
and quotient soundness.
\end{abstract}

\keywords{combinatory logic, combinator $S$, infinitary normal form, convertibility, head normal form}
\maketitle

\section{Introduction}

$\Ter$ is the set of closed terms built from $S$ by application, with $S\,x\,y\,z\to x\,z\,(y\,z)$
at any position; $\red$ is the reflexive--transitive closure of $\to$, and $\conv$ is
convertibility. The \emph{size} $\sz t$ is the number of occurrences of $S$ in $t$. Every
$t\in\Ter$ is $S\,u_1\cdots u_k$ for a unique $k\ge0$; the $u_i$ are its \emph{items}, and $t$ is a
\emph{head normal form} (hnf) if $k\le2$. If $k\ge3$, the \emph{head step} is
$t\hstep u_1\,u_3\,(u_2\,u_3)\,u_4\cdots u_k$, and $\hred$ is its reflexive--transitive closure.
The head run of $t$ stops at its first hnf, \emph{the} hnf of $t$, if there is one; its number of
items is the \emph{kind} of $t$. A term is \emph{literal} if it is its own hnf. Appending
arguments does not change a head step, so $t\hred t'$ implies $t\,r\hred t'\,r$. We write $=$ for
syntactic identity and $x\,\Sk k$ for $x\,S\cdots S$ with $k\ge0$ arguments $S$.

Every $t\in\Ter$ has an infinitary normal form $t^\infty$~\citep{s3}. We use it through its finite
prefixes $\NF(d,t)$ \citep[remark ``$t^\infty$ in Lean'']{s3}: $\NF(0,t)$ is a hole; $\NF(d+1,S)=S$; and
$\NF(d+1,t)=\NF(d,f)\cdot\NF(d,w)$ whenever $t\hred f\,w$ and the pair is \emph{stable}: $f$ has no
hnf, or $f\,w$ is an hnf. These prefixes exist and are unique, and reduction does not change
them~\citep{s3}. We write $t\approx u$, \emph{$t$ and $u$ have the same tree}, if
$\NF(d,t)=\NF(d,u)$ for all $d$; \citep{s3} identifies $\NF(d,t)$ with the depth-$d$ prefix of
$t^\infty$. By Church--Rosser, $t\conv u$ implies $t\approx u$. The converse,
\begin{center}
\emph{Lock~1:}\enskip $t\approx u$ implies $t\conv u$,
\end{center}
is the conjecture of~\citep{s6}, and by~\citep{s6} it would make convertibility of $S$-terms
decidable. It is open. We say that Lock~1 holds on the class of $t$ if every $u\approx t$ satisfies $u\conv t$.
It holds on the class of every normal term (Lemma~\ref{lem:nf}).

Put $e=S\,S$, $Q=S\,e\,e$, $G_0=Q$ and $G_{i+1}=e\,G_i$, so that $G_i=e^i\,Q$, and
$N(n,j)=G_j\,G_{n+1}$ for $0\le j\le n$. The $G_i$ are normal and pairwise distinct, since
$\sz{G_i}=5+2i$.

\begin{theorem}\label{thm:main}
Let $i\ge0$ and $m\ge1$. If $Y\in\Ter$ and $Y\,G_m\approx G_i\,G_m$, then $Y=G_i$.
\end{theorem}

\begin{corollary}\label{cor:lock}
Let $i\ge0$ and $m\ge1$. Every $u\in\Ter$ with $u\approx G_i\,G_m$ satisfies $u\conv G_i\,G_m$. In
particular, $u\approx N(n,j)$ implies $u\conv N(n,j)$.
\end{corollary}

The terms $G_i\,G_m$ do not normalize (Lemma~\ref{lem:family}), so Lemma~\ref{lem:nf} does not
apply to them. Their trees are infinite~\citep{s3}, and every node on the rightmost path has an hnf; this
separates them from the terms of~\citep{s11,s12}, which have none (Remark~\ref{rem:bot}).

\emph{Assumptions.} We use Church--Rosser~\citep{terese}; the prefix relation of~\citep{s3}, with
the existence, uniqueness and reduction invariance of its prefixes; the 38-element block model
of~\citep{s1}, whose fold $b$ is invariant under reduction and has $b(t)=36$ iff $t$ has no hnf;
the theorem of~\citep{s13} that $t\approx u$ implies $b(t)=b(u)$; and the finite facts about the
table of $b$ listed in Lemma~\ref{lem:table}, which are checked by kernel evaluation (Data
availability). Every case
split is on blocks, which are decidable, or on the syntax of a term, and the inductions are on $i$
and on size. The proofs are constructive.

\section{Trees and head runs}

\begin{lemma}\label{lem:K}
Let $t\approx u$. Then $t$ has an hnf iff $u$ has one. If the hnf of $t$ is $S\,a_1\cdots a_k$,
then the hnf of $u$ is $S\,a'_1\cdots a'_k$, with the same $k$ and $a_j\approx a'_j$.
\end{lemma}

\begin{proof}
If $t$ has the hnf $h$, then $t\red h$, so $\NF(d,t)=\NF(d,h)$. The pairs $(S,a)$ and $(S\,a,c)$
are stable, so $\NF(3,S)=S$, $\NF(3,S\,a)=S\cdot\NF(2,a)$ and
$\NF(3,S\,a\,c)=(S\cdot\NF(1,a))\cdot\NF(2,c)$. If $t$ has no hnf, every
stable pair $f\,w$ on its run has $f$ without an hnf, and by induction on $d$ the leftmost path of
$\NF(d,t)$ is $d$ application nodes ending in a hole. So $\NF(3,\cdot)$ determines the kind. Then
$\NF(d+1,S\,a)=S\cdot\NF(d,a)$ and $\NF(d+2,S\,a\,c)=(S\cdot\NF(d,a))\cdot\NF(d+1,c)$ give
$a_j\approx a_j'$.
\end{proof}

\begin{lemma}\label{lem:nf}
If $A$ is normal and $a\approx A$, then $a\red A$. Hence distinct normal terms do not have the same
tree.
\end{lemma}

\begin{proof}
Induction on $A$. A normal term has a head step unless it is an hnf, and its items are normal, so
$A$ is $S$, $S\,A_1$ or $S\,A_1\,A_2$ with $A_j$ normal. By Lemma~\ref{lem:K},
$a\hred S\,a_1\cdots a_k$ with the same $k$ and $a_j\approx A_j$, so $a_j\red A_j$ by induction,
and $a\red A$. A normal $a$ has no reduct but itself.
\end{proof}

\begin{lemma}\label{lem:head}
\textup{(a)} If $t\hred S\,a\,c$ in at least one step, the last step is $S\,S\,y\,z\hstep S\,z\,(y\,z)$,
so $a=z$ and $c=y\,z$.

\textup{(b)} If $t\hred S\,S\,c$, then $t=S\,S\,q\,\Sk k$ and $c=q\,\Sk k$ for some $q$ and $k\ge0$.
\end{lemma}

\begin{proof}
A head step from $S\,x\,y\,z\,r_1\cdots r_l$ gives a term whose items are those of $x$, followed by
$z$, $y\,z$, $r_1,\dots,r_l$. (a) Two items force $x=S$ and $l=0$. (b) Suppose the result is
$S\,S\,A\,\Sk k$. All its items but the second are $S$, and $y\,z$ is an application, so $y\,z$ is
the second item; hence $x$ has no items and $x=S$; then $z=S$, $A=y\,S$, and the step comes from $S\,S\,y\,\Sk{(k+1)}$.
Go back along the run from $S\,S\,c$, where $k=0$.
\end{proof}

\begin{lemma}\label{lem:family}
\textup{(a)} For all $i,m\ge0$,
\[
G_{i+1}\,G_m\hstep S\,G_m\,(G_i\,G_m),\qquad G_0\,G_m\hred S\,G_{m+1}\,(G_m\,G_{m+1}).
\]
So the hnf of $G_i\,G_m$ is $S\,A\,(\Gamma\,A)$, where $(\Gamma,A)=(G_{i-1},G_m)$ if $i\ge1$ and
$(\Gamma,A)=(G_m,G_{m+1})$ if $i=0$; and $\Gamma\,A=G_{i'}\,G_{m'}$ with $m'\ge m$.

\textup{(b)} No $G_i\,G_m$ normalizes.
\end{lemma}

\begin{proof}
(a) $G_{i+1}\,G_m=S\,S\,G_i\,G_m$, and
\[
Q\,G_m=S\,e\,e\,G_m\hstep S\,S\,G_m\,(e\,G_m)\hstep S\,(e\,G_m)\,(G_m\,(e\,G_m)).
\]
(b) If $t=G_i\,G_m\red N$ with $N$ normal, then $N\approx t$, so by Lemma~\ref{lem:K}
$N=S\,N_1\,N_2$ with $N_2\approx\Gamma\,A$, and Lemma~\ref{lem:nf} gives $\Gamma\,A\red N_2$. As
$\sz{N_2}<\sz N$ and $\Gamma\,A$ is again some $G_{i'}\,G_{m'}$, this cannot go on forever.
\end{proof}

\section{The table}

The fold is $b(S)=0$ and $b(f\,x)=\beta(b(f),b(x))$, where $\beta$ is the $38\times38$ table
of~\citep{s1}. As $b$ is invariant under reduction, an hnf $S$, $S\,a$ or $S\,a\,c$ of $t$ gives
$b(t)=0$, $\beta(0,b(a))$ or $\beta(\beta(0,b(a)),b(c))$.

\begin{lemma}\label{lem:table}
\textup{(T1)} No entry of $\beta$ is $0$, and $\beta(x,y)=1$, $2$ or $13$ only for $(x,y)=(0,0)$,
$(0,1)$ or $(2,1)$ respectively.
\textup{(T2)} $\beta(1,13)=\beta(1,12)=12$ and $\beta(13,12)=\beta(12,12)=25$.
\textup{(T3)} $\beta(y,12)=25$ iff $y\in\{11,12,13\}$.
\textup{(T4)} For all blocks $\alpha,\gamma$:
$\beta(\beta(0,\alpha),\gamma)\ne11$, and $\beta(\beta(0,\alpha),\gamma)=12$ only if $\alpha=0$.

\noindent Hence \textup{(a)} $b(u)=0$ iff $u=S$, and $b(u)=13$ iff $u=Q$; \textup{(b)}
$b(G_0)=13$, $b(G_i)=12$ for $i\ge1$, and $b(G_i\,G_m)=25$ for $i\ge0$, $m\ge1$.
\end{lemma}

\begin{proof}
(T1)--(T4) are finite computations on the table, checked by kernel evaluation (Data
availability). (a) By (T1), $b(u)=13$ forces
$u=f\,x$ with $b(f)=2$ and $b(x)=1$; then $x=S\,S$ and $f=S\,(S\,S)$, so $u=Q$. Conversely
$b(e)=\beta(0,0)=1$, $b(S\,e)=2$ and $b(Q)=13$. (b) By (a) and (T2), by induction on $i$.
\end{proof}

\section{Proof of Theorem~\ref{thm:main}}

By induction on $i$, for all $Y$. By~\citep{s13} and Lemma~\ref{lem:table}(b),
$\beta(b(Y),12)=b(Y\,G_m)=b(G_i\,G_m)=25$, so $b(Y)\in\{11,12,13\}$ by (T3). So $Y$ has an hnf $H$,
with $b(H)=b(Y)\ne0$, so $H\ne S$. Throughout we compare the hnf of $Y\,G_m$ with that of
$G_i\,G_m$ (Lemma~\ref{lem:family}(a)) by Lemma~\ref{lem:K}.

\emph{$H=S\,p$.} Then $Y\,G_m\hred S\,p\,G_m$, an hnf whose second item is $G_m$. The second item
of the hnf of $G_i\,G_m$ is $G_{i'}\,G_{m'}$ with $m'\ge1$. So $G_m\approx G_{i'}\,G_{m'}$, but
$b(G_m)=12$ and $b(G_{i'}\,G_{m'})=25$, against~\citep{s13}.

So $H$ has two items, and $b(Y)\ne11$ by (T4).

\emph{$b(Y)=13$.} Then $Y=Q=G_0$. If $i\ge1$, the first items of the hnfs of $Q\,G_m$ and
$G_i\,G_m$ are $G_{m+1}$ and $G_m$, against Lemma~\ref{lem:nf}.

\emph{$b(Y)=12$.} Then $H=S\,p\,q$ with $b(p)=0$ by (T4), so $p=S$, and
$Y\,G_m\hred S\,S\,q\,G_m\hstep S\,G_m\,(q\,G_m)$. If $i=0$, the first items are $G_m$ and
$G_{m+1}$, against Lemma~\ref{lem:nf}. If $i\ge1$, the second items give
$q\,G_m\approx G_{i-1}\,G_m$, so $q=G_{i-1}$ by induction. By Lemma~\ref{lem:head}(b),
$Y=S\,S\,q_0\,\Sk k$ with $q=q_0\,\Sk k$. The last argument of $G_{i-1}$ is $G_{i-2}$ or, for
$i=1$, $e$; neither is $S$, so $k=0$ and $Y=S\,S\,G_{i-1}=G_i$. \qed

\begin{remark}
With the theorem of~\citep{s13} that $\approx$ is a congruence, $Y\approx G_i$ gives
$Y\,G_1\approx G_i\,G_1$, so $Y=G_i$: the term $G_i$ is the only member of its class.
\end{remark}

\section{Proof of Corollary~\ref{cor:lock}}

By induction on $\sz u$, for all $i$ and $m\ge1$ at once. Let $t=G_i\,G_m$, with hnf
$S\,A\,(\Gamma\,A)$ as in Lemma~\ref{lem:family}(a). By Lemmas~\ref{lem:K} and~\ref{lem:nf},
$u\hred S\,a\,c$ with $a\red A$ and $c\approx\Gamma\,A=G_{i'}\,G_{m'}$, $m'\ge1$.

If $u$ is literal, $c$ is a proper subterm of $u$, so $c\conv\Gamma\,A$ by induction, and
$u=S\,a\,c\conv S\,A\,(\Gamma\,A)\conv t$.

Otherwise $a=z$ and $c=y\,z$ by Lemma~\ref{lem:head}(a). As $z\red A$, we get
$y\,A\approx y\,z\approx\Gamma\,A$, where $A=G_m$ or $G_{m+1}$ and $\Gamma=G_{i-1}$ or $G_m$. By
Theorem~\ref{thm:main}, $y=\Gamma$. So $u\hred S\,z\,(\Gamma\,z)\red S\,A\,(\Gamma\,A)$, the hnf
of $t$, and $u\conv t$. The terms $N(n,j)$ are $G_j\,G_{n+1}$. \qed

\section{Scope}\label{sec:scope}

\begin{remark}[Other families]\label{rem:families}
Let $c=S\,(S\,S)$, $D=S\,c\,c$, $X=S\,c\,e$, $H_j=c^j\,D$ and $E_j=e^j\,X$. Since
$c\,y\,z\hred S\,(y\,z)\,(z\,(y\,z))$ and $e\,y\,z\hstep S\,z\,(y\,z)$, the trees of
$H_j\,H_{n+1}$ and $E_j\,E_{n+1}$ also have an infinite path of nodes with an hnf. The proof of
Corollary~\ref{cor:lock} does not carry over. For $j\ge1$,
$H_j\,W\hred S\,(H_{j-1}\,W)\,(W\,(H_{j-1}\,W))$ with $W=H_{n+1}$, so the path passes through the
\emph{first} item, here $H_{j-1}\,H_{n+1}$; for $j=0$ it is $H_{n+1}\,H_{n+2}$. The argument of
Lemma~\ref{lem:family}(b), applied to first items, shows that these terms do not normalize. At a node $u\hred S\,z\,(y\,z)$ reached by a step,
Lemma~\ref{lem:nf} no longer forces $z$, and $z$ need not be a subterm of $u$, so the induction on
size stops. The terms $E_j\,E_{n+1}$ with $j\ge1$ have the hnfs of the $G$-family, but
$X\,W\hred S\,N\,(e\,W\,N)$ with $W=E_{n+1}$ and $N=W\,(e\,W)=E_{n+1}\,E_{n+2}$, so the first
item is again a term of the family, without a normal form. We do not know whether Lock~1 holds on these classes.
\end{remark}

\begin{remark}[Nodes without an hnf]\label{rem:bot}
For $R$ without an hnf, let $\tau_n(R)$ be its teeth~\citep{s11}. The prefix rules give that
$u\approx R$ iff $u$ has no hnf and $\tau_n(u)\approx\tau_n(R)$ for all $n\ge1$. (Sketch: the first
stable pairs of $u$ and $R$ have parts with the same tree, by uniqueness of prefixes, and the
converse is an induction on $d$.) Say that CT holds at $R$ if every $q$ without an hnf with
$\tau_n(q)\conv\tau_n(R)$ for all $n$ satisfies $q\conv R$. Granted the sketch, Lock~1 on the class
of $R$ implies CT at $R$; and if every tooth $\tau_n(R)$ satisfies Lock~1 on its class, CT at $R$
implies Lock~1 on the class of $R$. So any argument that settles Lock~1 at such a node proves CT
there. With Lemma~\ref{lem:nf}, the
results of~\citep{s12} for terms all of whose teeth normalize, such as $U_j\,U_m$ there, give
Lock~1 on their classes, under the same sketch. Our method treats only nodes with an hnf and adds
nothing at nodes without one.
\end{remark}

\begin{remark}[Finite models]
Lock~1 implies that the fold of $\Ter$ into every finite model of the $S$-law is constant on
$\approx$-classes. Corollary~\ref{cor:lock} shows this only on the classes of the $G_i\,G_m$; a term
whose values in a finite model follow those of $G_i\,G_m$ along the path need not have the tree of
$G_i\,G_m$. So the corollary proves this invariance for no finite model for which it was not
already known.
\end{remark}

\section*{Data availability}
\sloppy
Theorem~\ref{thm:main}, Corollary~\ref{cor:lock}, the Remark after Theorem~\ref{thm:main}, and
Lemmas~\ref{lem:K}--\ref{lem:table} except Lemma~\ref{lem:family}(b) are mechanized in
\texttt{lean-ski}~\citep{leanski} (Lean~4~\citep{lean4}), version~1.24.0 (merge commit
\texttt{8874b20}, released 27 November 2026), module \texttt{Ski/NoImp.lean}. They use only the
axioms \texttt{propext} and \texttt{Quot.sound}, and \texttt{Ski/Audit.lean} checks every
declaration of the module. The Lean names are:
\begin{itemize}
\item Theorem~\ref{thm:main}: \texttt{NoImp.noimp}; the Remark after it: \texttt{NoImp.eqv\_G}.
\item Corollary~\ref{cor:lock}: \texttt{NoImp.lock1\_G}, which gives a common reduct of $u$ and
  $G_i\,G_m$; \texttt{NoImp.lock1\_G\_conv}, convertibility; and \texttt{NoImp.lock1\_N}, for
  $N(n,j)$ with any $n$ and $j$. They prove the stated direction, that $u\approx G_i\,G_m$ implies
  $u\conv G_i\,G_m$.
\item Lemma~\ref{lem:K}: a form of it is \texttt{Congruence.factK}. Lemma~\ref{lem:nf}:
  \texttt{NoImp.nf\_reduces}; its second sentence is mechanized for the $G_i$ only
  (\texttt{NoImp.G\_eqv\_inj}), which is the case used.
\item Lemma~\ref{lem:head}: (a) \texttt{NoImp.last\_step}, (b) \texttt{NoImp.comb}.
  Lemma~\ref{lem:family}(a): \texttt{NoImp.G\_hnf}.
\item Lemma~\ref{lem:table}: (T1) is \texttt{FFBlock.eq\_S\_of\_fold} (no entry is $0$) and
  \texttt{NoImp.pre1}, \texttt{pre2}, \texttt{pre13}; (T3) is \texttt{NoImp.B3}; (T4) is
  \texttt{NoImp.B4\_two}; (T2) is evaluated inside \texttt{NoImp.fold\_G\_succ} and
  \texttt{NoImp.fold\_GG}, which are~(b); and (a) is \texttt{NoImp.rig13}. The table facts are
  proved by \texttt{decide +kernel} on \texttt{blockTable} in \texttt{Ski/HnfS.lean}.
\end{itemize}
Not mechanized: Lemma~\ref{lem:family}(b), the remarks of Section~\ref{sec:scope}, including the
sketch in Remark~\ref{rem:bot}, and the identification of the prefixes $\NF(d,t)$ with prefixes
of $t^\infty$, which no proof here uses. The inputs are mechanized in the same release: the prefix
relation is \texttt{InfNF.NF}, with \texttt{nf\_exists}, \texttt{nf\_unique} and
\texttt{nf\_wsteps}; Church--Rosser is \texttt{church\_rosser}; the invariance of the block under
$\approx$ is \texttt{FFBlock.fold\_eqv}; and the congruence used in the Remark after
Theorem~\ref{thm:main} is \texttt{Congruence.eqv\_app}. In \texttt{lean-ski}, $\approx$ is
\texttt{Congruence.Eqv}, the relation of this paper. Convertibility there is weak convertibility
over $SKI$; since $S$-terms reduce only to $S$-terms, by Church--Rosser it agrees with $\conv$ on
$S$-terms (\texttt{Conv.conv\_iff\_join}).

\bibliographystyle{plainnat}
\bibliography{refs}
\end{document}
